% !TEX program = xelatex
% =====================================================================
%  TUBITAK 2204-A  --  Proje Detay Raporu
%  Sosyal Medya Duygu Zaman Serilerinin Çoklu Fraktal Analizi
%  Derleyici: XeLaTeX
% =====================================================================
\documentclass[11pt,a4paper]{article}

\usepackage{fontspec}
\defaultfontfeatures{Ligatures=TeX}
\usepackage{polyglossia}
\setmainlanguage{turkish}
\setotherlanguage{english}

\usepackage{amsmath,amssymb}
\usepackage{booktabs}
\usepackage{longtable}
\usepackage{array}
\usepackage{graphicx}
\usepackage{geometry}
\geometry{margin=2.4cm}
\usepackage{enumitem}
\usepackage{float}
\usepackage{caption}
\captionsetup{font=small,labelfont=bf}
\usepackage{hyperref}
\hypersetup{colorlinks=true,linkcolor=blue,urlcolor=blue}
\usepackage{xcolor}

\newcommand{\Dal}{\Delta\alpha}

\title{\bfseries Sosyal Medya Duygu Zaman Serilerinin Çoklu Fraktal Analizi ile\\
Kaotik ve Stokastik Dinamiklerin Ayrıştırılması}
\author{TÜBİTAK 2204-A Proje Detay Raporu\\[3pt]
\small Matematik Alanı --- Algoritma Tasarımı ve Uygulamaları}
\date{\today}

\begin{document}
\maketitle

\begin{center}
\begin{minipage}{0.92\textwidth}
\small\itshape
Bu rapor, konuya tamamen yabancı bir okuyucunun (jüri üyesi, öğretmen, öğrenci)
her adımı takip edebilmesi için yazılmıştır. Bu nedenle teknik kavramlar önce
gündelik benzetmelerle sezdirilir, sonra matematiksel olarak tanımlanır.
Rapor tek bir soruya odaklanır: \textbf{Bir deprem sırasında milyonlarca
insanın sosyal medyada birlikte ürettiği duygu zaman serisi tamamen rastgele
(stokastik) midir, yoksa kendine özgü, ölçekten bağımsız bir düzeni (çoklu
fraktal yapı) taşır mı?}
\end{minipage}
\end{center}

\vspace{0.5cm}

% =====================================================================
\section*{Özet}
\addcontentsline{toc}{section}{Özet}
% =====================================================================

Bu projede, büyük bir toplumsal sarsıntı (deprem) sırasında sosyal medyada
oluşan \emph{kolektif duygunun} bir zaman serisi olarak nasıl davrandığı
incelenmiştir. İki bağımsız deprem olayında (6 Şubat 2023 Kahramanmaraş,
30 Ekim 2020 İzmir) toplam yaklaşık 1{,}9 milyon Türkçe tweet toplanmış;
her tweet, altı temel duygu boyutunda (öfke, tiksinme, korku, sevinç, üzüntü,
şaşkınlık) puanlanmış ve dakika bazında ortalama alınarak sürekli duygu zaman
serileri üretilmiştir. Bu serilerin istatistiksel doğası \emph{çoklu fraktal
dalgalanma analizi} (MF--DFA) ile ölçülmüş, gözlemin rastgelelikten gerçekten
ayrışıp ayrışmadığı \emph{surrogate (vekil veri) testleri} ile sınanmıştır.

Sonuçlar üç başlıkta özetlenebilir. Birincisi, duygu zaman serisi \emph{tamamen
rastgele değildir}; gözlenen çoklu fraktallik genişliği ($\Dal$), aynı dağılıma
ve aynı doğrusal korelasyona sahip vekil verilerin ürettiğinden istatistiksel
olarak anlamlı biçimde büyüktür ($z\approx +10$--$+22$). Yani duyguda,
doğrusal korelasyonla açıklanamayan, gerçek bir \emph{doğrusal olmayan} yapı
vardır. İkincisi, bu yapı \emph{zamanla değişir}: kriz anında (depremden
sonraki ilk saatler) çoklu fraktallik güçlüdür; günler geçip toplum normale
döndükçe zayıflar ($\Delta\alpha$ Kahramanmaraş'ta $0.78\to0.38$, İzmir'de
$0.67\to0.25$). Üçüncüsü, bulgu iki bağımsız olayda, iki farklı duygu modeliyle
ve farklı veri hacimleriyle tekrarlanmıştır; dolayısıyla tek bir olaya özgü
tesadüf değildir.

Projenin bir diğer önemli çıktısı metodolojiktir: sosyal medya duygu ölçümünde
naif bir model seçiminin sistematik ve ciddi hata üretebileceği somut
örneklerle gösterilmiş ve bu hata ayıklanarak giderilmiştir. Ayrıca
surrogate testlerinin kullanımına ilişkin bir sınırlamanın varlığı
deneysel olarak belirlenmiş ve belgelenmiştir: IAAFT null'u, \emph{ayrık} genlik
dağılımlı serilerde çoklu fraktalliği yok edemeyebilmekte; bu nedenle tüm
sonuçlar iki bağımsız null modeliyle (IAAFT ve shuffle) birlikte raporlanmıştır.

\textbf{Anahtar kelimeler:} çoklu fraktal analiz, MF--DFA, Hurst üssü,
surrogate testi, IAAFT, duygu zaman serisi, sosyal medya, deprem,
doğrusal olmayan dinamikler, uzun erimli bellek.

\tableofcontents
\newpage

% =====================================================================
\section{Giriş}
% =====================================================================

\subsection{Bir soruyla başlayalım}
% ---------------------------------------------------------------------

Bir şehirde büyük bir deprem olduğunu düşünün. Dakikalar içinde milyonlarca
insan telefonuna sarılır, sosyal medyada yazar, paylaşır, ağlar, yardım ister,
yardım eder. Bu insanlar birbirini tanımaz, birbirleriyle konuşmaz; ama hep
birlikte, devasa bir ``duygu dalgası'' üretirler. Peki bu dalga \emph{nasıl bir
şeydir}?

Bu dalgayı ölçmenin basit bir yolu vardır. Zamanı birer dakikalık aralıklara
böleriz. Her dakikada, o dakika içinde atılan tüm tweet'leri alıp ``bu tweet'ler
ne kadar üzüntü içeriyor?'' sorusunu sorar ve tek bir sayı elde ederiz. Böylece
elimizde bir sayı dizisi olur:

\begin{center}
$0.42,\ 0.44,\ 0.41,\ 0.55,\ 0.61,\ 0.58,\ 0.43,\ \dots$
\end{center}

İşte bu dizinin adı \textbf{zaman serisidir}: ardışık zamanlarda ölçülmüş
değerlerin sıralı listesi. Şimdi asıl soru şudur:

\begin{quote}
\emph{Bu sayılar arasında bir düzen var mıdır, yoksa sadece rastgele mi
sıralanmışlardır?}
\end{quote}

Eğer sayılar gerçekten rastgeleyse (her dakika bir zar atılmış gibi), o zaman
geçmişe bakarak geleceği tahmin etmek temelde imkânsızdır ve toplumun duygusu
hakkında ``yapısal'' hiçbir şey söyleyemeyiz. Ama eğer sayılar arasında
zamana yayılmış, \emph{ölçekten bağımsız} bir düzen varsa, o zaman ortada
ciddi bir olgu vardır: insanlar farkında olmadan ortak bir karmaşık dinamiğe
girmektedirler. Bu proje tam olarak bu ikisini \emph{birbirinden ayırmayı}
amaçlar.

\subsection{Neden ``rastgele mi, fraktal mi?'' sorusu önemlidir}
% ---------------------------------------------------------------------

Bilimde bir olayı anlamanın ilk adımı, onun \emph{rastgele} mi yoksa
\emph{belirli bir yapı taşıyor mu} olduğunu bilmektir. Bu ayrım sadece akademik
bir merak değildir; pratik sonuçları vardır:

\begin{itemize}[leftmargin=1.4em]
  \item Eğer duygu serisi rastgeleyse, onu modellemek için yalnızca
        ``ortalama'' ve ``dağılım'' yeterlidir; başka bir yapı aramaya gerek
        yoktur.
  \item Eğer seri \emph{uzun erimli bellek} (bugünkü değerin, günler önceki
        değerden hâlâ etkilenmesi) ve \emph{çoklu fraktallik} (farklı
        büyüklükteki dalgalanmaların farklı ölçeklerde farklı davranması)
        taşıyorsa, o zaman:
        \begin{itemize}
          \item topyekûn duygunun ``hafızası'' vardır; kriz yönetimi,
                iletişim ve psikolojik destek planlamasında bu hafıza
                kullanılabilir;
          \item kriz ile normal dönem istatistiksel olarak birbirinden
                ayrılabilir; bu, bir \emph{erken uyarı göstergesi} (early
                warning indicator) fikrinin temelidir;
          \item ``insan toplulukları, tıpkı türbülans ya da finansal piyasalar
                gibi, doğrusal olmayan (nonlinear) dinamikler üretir'' tezi
                somut veriyle sınanmış olur.
        \end{itemize}
\end{itemize}

Kısacası, bu proje bir \emph{ayrıştırma} problemidir: karmaşık görünen bir
sinyalin içindeki ``rastgele'' ile ``yapılı'' kısmı birbirinden ayırmak.

\subsection{Neden sosyal medya? Neden deprem?}
% ---------------------------------------------------------------------

\textbf{Neden sosyal medya?} Çünkü insan davranışını, üstelik \emph{zaman
damgalı} ve \emph{ölçeklenebilir} biçimde gözlemlemenin nadir yollarından
biridir. Klasik anketler geçmişe dönük ve seyrektir; oysa sosyal medya,
saniyeler mertebesinde, doğal ortamda ve sürekli veri üretir. Bir şehrin
``toplu ruh hâli''ni dakika dakika izleyebileceğimiz bir kayıttır.

\textbf{Neden deprem?} Çünkü deprem, tetikleyicisi (yer sarsıntısı) kesin
olarak bilinen, ani ve çok şiddetli bir toplumsal olaydır. Bu sayede
``kriz başladı'' anını saniyesi saniyesine bilebiliriz ki bu, bilimsel olarak
paha biçilmezdir: kriz öncesi, kriz anı ve kriz sonrasını net biçimde
ayırabiliriz. Ayrıca depremler, birbirinden bağımsız olarak gerçekleştiği için
elde edilen sonucun \emph{tekrarlanabilirliğini} test etmemize olanak tanır.

\subsection{Literatür ve boşluk}
% ---------------------------------------------------------------------

Fraktal ve çoklu fraktal analiz, 1990'lardan bu yana pek çok alanda
kullanılmıştır:
\begin{itemize}[leftmargin=1.4em]
  \item \textbf{Fizik:} türbülans, akışkan dinamiği, dalga kırılması.
  \item \textbf{Finans:} borsa getirileri, döviz kurları, volatilite.
  \item \textbf{Biyoloji/tıp:} kalp atışı aralıkları (HRV), beyin
        aktivitesi (EEG), DNA dizileri.
  \item \textbf{İklim:} sıcaklık, yağış, nehir akışı.
  \item \textbf{Karmaşık sistemler:} internet trafiği, trafik akışı.
\end{itemize}

Buna karşılık, \emph{sosyal medya duygu zaman serilerinin} çoklu fraktal
analizi çok yenidir ve özellikle \emph{afet bağlamında} neredeyse bakir bir
alandır. Mevcut çalışmaların çoğu ya duygu sınıflandırmasıyla (makine
öğrenmesi) ya da tek ölçekli istatistiklerle (örneğin basit otokorelasyon)
sınırlı kalmıştır. Oysa kolektif duygunun doğasını anlamak için
\emph{ölçekten bağımsız} bir çerçeve gerekir. Bu proje tam olarak bu boşluğu
hedefler: \emph{iki bağımsız afet olayında, duygu zaman serisini çoklu fraktal
analizle incelemek ve rastgelelikle sistematik olarak karşılaştırmak.}

\subsection{Amaç, araştırma soruları ve hipotezler}
% ---------------------------------------------------------------------

\textbf{Genel amaç:} Sosyal medyada bir kriz anında oluşan kolektif duygu
zaman serisinin istatistiksel doğasını, rastgele (stokastik) bir süreçten
ayırarak ortaya koymak.

\textbf{Araştırma soruları:}
\begin{enumerate}[leftmargin=1.6em,label=\textbf{AS\arabic*.}]
  \item Duygu zaman serisi, beyaz gürültüye (tamamen rastgele) kıyasla
        anlamlı bir çoklu fraktal genişlik ($\Dal$) sergiliyor mu?
  \item Bu genişlik, aynı dağılımı ve aynı \emph{doğrusal} korelasyonu koruyan
        vekil verilere (IAAFT) kıyasla da anlamlı mı? (Yani doğrusal olmayan
        bir yapı var mı?)
  \item Çoklu fraktallik \emph{kriz} ve \emph{normalleşme} fazları arasında
        değişiyor mu?
  \item Duygunun karmaşıklığı, salt \emph{aktivite yoğunluğundan} (tweet
        hacmi) ayrılabiliyor mu?
  \item Bu sonuçlar iki bağımsız deprem olayında tekrarlanıyor mu
        (genellenebilirlik)?
\end{enumerate}

\textbf{Hipotezler:}
\begin{description}[leftmargin=1.6em]
  \item[H1] Duygu serisi, rastgele vekillere kıyasla anlamlı $\Dal$ gösterir
        (yani rastgele değildir).
  \item[H2] Bu çoklu fraktallik, doğrusal korelasyonla açıklanamaz
        (IAAFT testinde de anlamlı kalır); dolayısıyla doğrusal olmayan
        dinamiklerin izidir.
  \item[H3] Kriz fazı, normalleşme fazından daha güçlü çoklu fraktal yapıya
        sahiptir.
  \item[H4] Duygu, tweet hacminden istatistiksel olarak ayrılabilir.
        \emph{(Bu hipotez yalnızca kısmen desteklenmiştir; bkz.
        Bölüm~\ref{sec:sinirlilik}.)}
\end{description}

\subsection{Özgün katkılar}
% ---------------------------------------------------------------------

\begin{enumerate}[leftmargin=1.6em]
  \item \textbf{Konu seçimi:} Çoklu fraktal analizin afet sonrası kolektif
        duyguya uygulanması özgündür.
  \item \textbf{İki olay tasarımı:} Farklı büyüklük ve farklı duygusal
        profildeki iki bağımsız depremle genellenebilirlik test edilmiştir.
  \item \textbf{Metodolojik titizlik:} Duygu ölçümünde model seçiminin
        yanlış sonuç üretebileceği somut gösterilerek düzeltilmiş; analiz
        parametreleri (ölçek aralığı, pencere) bağımsız kriterlerle
        (örneğin $R^2$ doğrusallık ölçütü) sabitlenmiştir.
  \item \textbf{Doğrusal/doğrusal-olmayan ayrımı:} Surrogate hiyerarşisi
        (shuffle, faz-rastgele, IAAFT) ile ``sadece korelasyon mu, yoksa
        doğrusal olmayan yapı mı?'' sorusu net biçimde ayrıştırılmıştır.
  \item \textbf{Reprodüklenebilirlik:} Tüm analiz açık kaynak kodla,
        sıfırdan yazılan bir MF--DFA motoruyla yapılmış ve bilinen referans
        süreçlerle (p-model, kesirli Gauss gürültüsü, multifraktal rastgele
        yürüyüş) doğrulanmıştır. 28 adet birim testi içeren otomatik test
        paketi, matematiksel doğruluğun sürekli denetlenmesini sağlar.
  \item \textbf{Null model sınırlamasının belgelenmesi:} IAAFT null'unun
        ayrık genlik dağılımlı serilerde çoklu fraktalliği yok edemediği
        deneysel olarak gösterilmiş ve bu durumda çift null testi
        (IAAFT + shuffle) kullanılması önerilmiştir.
\end{enumerate}

% =====================================================================
\section{Kuramsal Arka Plan: Sıfırdan Anlatım}
% =====================================================================

Bu bölüm, raporun kalbidir. Burada kullanılan her kavram, hiçbir ön bilgi
varsayılmadan, önce sezgisel olarak, sonra matematiksel olarak anlatılır.

\subsection{Zaman serisi nedir?}
% ---------------------------------------------------------------------

Bir \textbf{zaman serisi}, belirli bir olayın zaman içinde sırayla ölçülmüş
değerlerinin listesidir. Gündelik örnekler:
\begin{itemize}[leftmargin=1.4em]
  \item Her gün aynı saatte ölçülen hava sıcaklığı;
  \item Her saniye kaydedilen borsa fiyatı;
  \item Her dakika hesaplanan ortalama duygu puanı.
\end{itemize}

Kritik nokta şudur: zaman serisinde \emph{sıra önemlidir}. Aynı sayıları
karıştırıp rastgele sırada dizsek, artık elimizde aynı \emph{küme} vardır ama
farklı bir \emph{seri} vardır. Az sonra göreceğimiz gibi, bu projenin en
güçlü aracı tam da bu ayrımdan yararlanır: ``Sayıların sırası bilgi taşıyor mu?''

Duygu serimizi matematiksel olarak $x_1, x_2, \dots, x_N$ biçiminde yazarız.
Burada $x_i$, $i$-inci dakikadaki ortalama duygu puanıdır ve $N$ toplam
dakika sayısıdır.

\subsection{Rastgelelik (stokastik) ve ``null model''}
% ---------------------------------------------------------------------

Bir sürecin \textbf{rastgele (stokastik)} olması, geleceğinin geçmişten
\emph{bağımsız} olması demektir. En saf hâli \textbf{beyaz gürültüdür}:
her değer, aynı dağılımdan bağımsız olarak çekilmiş bir zar atışı gibidir.
Beyaz gürültüde ``bugün büyükse yarın da büyük olur'' gibi bir eğilim yoktur.

Bilimde bir iddiayı test etmenin en sağlam yolu, onu bir \emph{zıt} iddiayla
yarıştırmaktır. Zıt iddiaya \textbf{null model} (sıfır hipotezi) denir.
Bizim durumumuzda doğal null model: \emph{``Duygu serisi, gözlenen dağılıma
sahip rastgele bir gürültüdür.''} Eğer gözlediğimiz yapı, bu null modelin
üretebileceğinden \emph{belirgin biçimde} farklıysa, rastgelelik iddiası
reddedilir.

Burada bir incelik vardır ve projemizin temel fikri buradadır: ``rastgele''
tek bir şey değildir. İki farklı derecesi vardır:
\begin{enumerate}[leftmargin=1.6em]
  \item \textbf{Tamamen rastgele} (beyaz gürültü): ne sıra ilişkisi ne
        yapı vardır.
  \item \textbf{Doğrusal korelasyonlu rastgele:} değerler rastgeledir ama
        \emph{birbirine bağlıdır} --- bugünkü değer, dünkü değere doğrusal
        olarak bağlıdır (örneğin $x_i = 0.9\,x_{i-1} + \text{gürültü}$).
        Bu tür süreçler ``hafızalı'' görünür ama \emph{doğrusal} kurallarla
        üretilebilir.
\end{enumerate}

Bu projenin özü, gözlenen karmaşıklığın \emph{sadece ikinci türden mi} (yani
doğrusal korelasyon) yoksa \emph{gerçekten doğrusal olmayan bir yapıdan mı}
kaynaklandığını ayırmaktır. Bunu yapabilmek için sayıların sırasını bozan ama
doğrusal özellikleri koruyan ``vekil veriler'' üreteceğiz (Bölüm~3.5).

\subsection{Uzun erimli bellek ve Hurst üssü}
% ---------------------------------------------------------------------

Bir seride ``hafıza'' olup olmadığını ölçmenin klasik yolu
\textbf{otokorelasyondur}: serinin kendisinin, kendisinin kaydırılmış hâliyle
ne kadar benzeştiği. Eğer komşu değerler birbirine benziyorsa pozitif
otokorelasyon vardır.

Bazı süreçlerde bu benzerlik \emph{çok yavaş} azalır; saatler, hatta günler
sonra bile sürer. Buna \textbf{uzun erimli bellek} denir. Böyle bir serinin
``hafızası''nı tek bir sayıyla özetleyebiliriz: \textbf{Hurst üssü} $H$.

Sezgisel anlamı:
\begin{itemize}[leftmargin=1.4em]
  \item $H = 0.5$: Bellek yok (saf rastgele yürüyüşün artışları gibi).
  \item $H > 0.5$: \emph{Kalıcı} (persistent). Bir yöne giden seri, o yönde
        gitmeye devam etme eğilimindedir. Örnek: bir nehrin akışı; uzun süre
        yüksek kalan bir taşkın.
  \item $H < 0.5$: \emph{Anti-kalıcı} (anti-persistent). Yüksek değeri düşük,
        düşük değeri yüksek izler; seri sürekli yön değiştirir. Örnek: bir
        kumarbazın ``şansının dönmesi''.
\end{itemize}

Örneğin, duygu serimizde $H \approx 0.8$ bulursak, bu şu demektir: bir kez
toplum ``kaygı moduna'' girdiğinde, bu hâl kendiliğinden hızla dağılmaz;
günler boyunca izini sürdürür. Bu, afet sonrası toplumsal iyileşme
süreçleriyle doğrudan ilgilidir.

\subsection{Fraktal ve ölçekten bağımsızlık}
% ---------------------------------------------------------------------

\textbf{Fraktal}, ``kendine benzer'' demektir: bir nesnenin parçası, bütününe
benzer; ne kadar yakından bakarsanız bakın, aynı türde yapı görürsünüz.
Klasik örnek: \emph{kıyı şeridi}. Bir kıyı çizgisini ne kadar hassas ölçerseniz
(1 km'lik cetvel yerine 1 m'lik, sonra 1 cm'lik), ölçtüğünüz uzunluk o kadar
büyür; kesin bir değere yakınsamaz. İşte bu ``ölçekten bağımsızlık,'' fraktal
geometrinin özüdür (Mandelbrot, 1982).

Zaman serilerinde de benzer bir kavram vardır. Bir serinin ``pürüzlülüğüne''
farklı zaman ölçeklerinden bakalım: 1 dakikaya, 1 saate, 1 güne. Eğer bu
pürüzlülük her ölçekte \emph{aynı kuralı} izliyorsa, seri fraktaldır. Matematiksel
ifadesi bir güç yasasıdır:
\begin{equation}
  F(s) \sim s^{H},
  \label{eq:powerlaw}
\end{equation}
burada $s$ baktığımız pencere uzunluğu (ölçek), $F(s)$ o ölçekteki
dalgalanmanın büyüklüğü, $H$ ise ölçekleme üssüdür. \eqref{eq:powerlaw} gibi
bir bağıntıya ``ölçekleme yasası'' denir; logaritmik eksende çizildiğinde
bir doğru verir, eğimi $H$'dir.

\subsection{Tekli (mono) ve çoklu (multi) fraktallik}
% ---------------------------------------------------------------------

Yukarıdaki $H$ \emph{tek} bir sayıdır. Peki ya bir seri, farklı büyüklükteki
dalgalanmalar için \emph{farklı} ölçekleme gösteriyorsa? Yani küçük
dalgalanmalar başka bir üsle, büyük dalgalanmalar başka bir üsle büyüyorsa?

Bu durumda tek bir $H$ yetmez. Seri \textbf{çoklu fraktaldır}. Çoklu
fraktallik, genellikle ``\emph{intermittency}'' (kesikli patlamalar) ile
birlikte görülür: uzun sakin dönemler, aralarda ani ve şiddetli patlamalar.

Gündelik benzetmeler:
\begin{itemize}[leftmargin=1.4em]
  \item Trafik akışı: çoğu zaman sakindir, ama aniden tıkanabilir.
  \item Borsa: uzun ``sakin'' dönemler, aralarda volatilite patlamaları.
  \item Türbülans: aynı akış içinde hem düzenli hem kaotik bölgeler.
\end{itemize}

Duygu serisinde bunun karşılığı açıktır: yeni bir şok haberi geldiğinde
(sözgelimi bir enkazdan çocuk çıkarıldığında veya bir artçı sarsıntı olduğunda)
duygu aniden patlar; sonra tekrar sakinleşir. İşte bu ``sakin \emph{ve}
patlamalı'' dokuyu ölçmenin aracı çoklu fraktal analizdir.

\subsection{DFA: Dalgalanma Analizi, adım adım}
% ---------------------------------------------------------------------

Çoklu fraktal analizin temeli, Peng ve arkadaşlarının (1994) geliştirdiği
\textbf{DFA} (Detrended Fluctuation Analysis, ``eğilimden arındırılmış
dalgalanma analizi'') yöntemidir. Adım adım anlatalım.

\textbf{1. Adım --- Ortalamayı çıkar.}
Serinin ortalamasını $\bar{x} = \frac{1}{N}\sum_{i=1}^{N} x_i$ hesaplar ve her
değerden çıkarırız. Bu, serinin etrafında ``sıfır etrafında dalgalandığı''
bir hâle gelmesini sağlar.

\textbf{2. Adım --- Profil oluştur (kümülatif toplam).}
\begin{equation}
  Y(i) = \sum_{k=1}^{i}\bigl(x_k - \bar{x}\bigr), \qquad i=1,\dots,N.
  \label{eq:profile}
\end{equation}
Bu adım, gürültüyü bir ``rastgele yürüyüşe'' çevirir; trendler ve uzun
dalgalanmalar böylece belirginleşir. $Y(i)$, artık bir yürüyüşün konumudur.

\textbf{3. Adım --- Pencerelere böl.}
Profili $Y(i)$'yi (Denklem~\eqref{eq:profile}), uzunluğu $s$ olan eşit parçalara
ayırırız. Toplam pencere sayısı $N_s = \lfloor N/s \rfloor$. (Uçta artan kısım
için aynı işlem sondan başa tekrarlanır ki hiçbir veri boşa gitmesin.)

\textbf{4. Adım --- Yerel trendi çıkar (``detrend'').}
Her pencerede, $Y(i)$'ye bir polinom (genellikle düz çizgi ya da parabol)
uyarlar ve \emph{çıkarırız}. Kalan ``artık'' yalnızca o penceredeki dalgalanmadır.
Bu, seriyi yerel trendlerden arındırır; MF--DFA'nın durağan olmayan
(nonstationary) verilerle çalışabilmesinin sırrı budur.

\textbf{5. Adım --- Dalgalanmayı ölç.}
$\nu$-uncu penceredeki artıkların ortalama karesini alalım: $F^2(\nu,s)$.
Tüm pencereler üzerinden ortalamasını alınca o ölçekteki dalgalanma:
\begin{equation}
  F(s) = \left[\frac{1}{N_s}\sum_{\nu=1}^{N_s} F^2(\nu,s)\right]^{1/2}.
  \label{eq:F2}
\end{equation}

\textbf{6. Adım --- Ölçekleme yasasını bul.}
Farklı $s$ değerleri için $F(s)$'i (Denklem~\eqref{eq:F2}) hesaplayıp
$\log F(s)$'i $\log s$'ye karşı çizeriz. Eğer bir doğru elde edersek, eğimi
bize Hurst üssü $H$'yi verir (Bkz. Denklem~\eqref{eq:powerlaw}). Doğrusallık
(yüksek $R^2$), ölçekleme yasasının geçerli olduğunun kanıtıdır; bu projede
bu ölçütü titizlikle kullandık (bkz. Bölüm~\ref{sec:olcek}).

\subsection{Çoklu Fraktal DFA (MF--DFA)}
% ---------------------------------------------------------------------

Klasik DFA, dalgalanmayı \emph{ortalama kare} (yani 2.\ mertebe) ile ölçer.
MF--DFA (Kantelhardt vd., 2002), bu tek ölçüyü bir \emph{ailesine}
genişletir: her pencereye bir $q$ üssü uygularız
(Denklem~\eqref{eq:Fq}),
\begin{equation}
  F_q(s) =
  \begin{cases}
    \left[\dfrac{1}{N_s}\displaystyle\sum_{\nu=1}^{N_s}
      \bigl(F^2(\nu,s)\bigr)^{q/2}\right]^{1/q}, & q \neq 0,\\[2.2ex]
    \exp\!\left[\dfrac{1}{2N_s}\displaystyle\sum_{\nu=1}^{N_s}
      \ln F^2(\nu,s)\right], & q = 0.
  \end{cases}
  \label{eq:Fq}
\end{equation}
$q$'nun anlamı şudur:
\begin{itemize}[leftmargin=1.4em]
  \item $q > 0$: \emph{büyük} dalgalanmalar baskın hâle gelir (patlamalara bakar).
  \item $q < 0$: \emph{küçük} dalgalanmalar baskın hâle gelir (sakin dönemlere bakar).
  \item $q = 0$: logaritmik ortalama (her ikisine de eşit ağırlık).
\end{itemize}

Sonra her $q$ için ayrı bir ölçekleme yasası ararız
(Denklem~\eqref{eq:hq}):
\begin{equation}
  F_q(s) \sim s^{\,h(q)}.
  \label{eq:hq}
\end{equation}
Buradaki $h(q)$ \textbf{genelleştirilmiş Hurst üssüdür}. Şimdi kilit ayrım
gelir:
\begin{itemize}[leftmargin=1.4em]
  \item Eğer $h(q)$ tüm $q$ için \emph{sabit} ise (yani $h(q)=H$), seri
        \emph{tekli fraktaldır}. Küçük ve büyük dalgalanmalar aynı kurala uyar.
  \item Eğer $h(q)$ $q$'ya göre \emph{değişiyorsa}, seri \emph{çoklu
        fraktaldır}. Farklı büyüklükteki dalgalanmalar farklı ölçekleme
        gösterir; bu, doğrusal olmayan yapının ve intermittency'nin işaretidir.
\end{itemize}

\subsection{Çoklu fraktal spektrum $f(\alpha)$ ve Legendre dönüşümü}
% ---------------------------------------------------------------------

$h(q)$ eğrisini, doğrudan yorumlaması biraz soyut olan $f(\alpha)$
spektrumuna çevirmek gelenekseldir. Önce ölçek üssü
(Denklem~\eqref{eq:tau}):
\begin{equation}
  \tau(q) = q\,h(q) - 1 .
  \label{eq:tau}
\end{equation}
$\tau(q)$, ``kütle üssü''dür. Ardından \textbf{Legendre dönüşümü} ile
\emph{singularite üssü} $\alpha$ ve spektrum yüksekliği $f(\alpha)$ elde edilir
(Denklem~\eqref{eq:legendre}):
\begin{equation}
  \alpha = \frac{d\tau}{dq}, \qquad f(\alpha) = q\,\alpha - \tau(q).
  \label{eq:legendre}
\end{equation}

Bunların anlamı:
\begin{itemize}[leftmargin=1.4em]
  \item $\alpha$, bir \emph{yerel pürüzlülük üssüdür} (Hölder üssü). Küçük
        $\alpha$ çok pürüzlü/patlamalı bölgeleri, büyük $\alpha$ düzgün
        bölgeleri temsil eder.
  \item $f(\alpha)$ ise ``$\alpha$ üssüne sahip noktalar kümesinin fraktal
        boyutudur''; yani o tür bölgelerin ne kadar ``yaygın'' olduğunu söyler.
  \item Spektrumun \emph{tepe noktası} ($f_{\max}$) genellikle $q=0$'a karşılık
        gelir ve boyut bilgisi taşır.
  \item \textbf{Spektrumun genişliği}
        \begin{equation}
          \Dal = \alpha_{\max} - \alpha_{\min}
          \label{eq:width}
        \end{equation}
        çoklu fraktalliğin \emph{şiddetidir}. Tekli fraktalda spektrum bir
        noktaya çöker ($\Dal \to 0$); çoklu fraktalda genişler.
\end{itemize}

Kısaca: \textbf{$\Dal$ (Denklem~\eqref{eq:width}) büyükse, seri hem sakin hem
patlamalı bölgeler barındırıyor ve bunlar farklı ölçekleme yasalarına tabi;
yani zengin, doğrusal olmayan bir dinamik vardır.}

\subsection{$\Dal$ tek başına yeter mi? Neden surrogate testleri?}
% ---------------------------------------------------------------------

Burada çok önemli bir tuzak vardır. $s<24$ dakika gibi \emph{çok küçük}
ölçeklerde ölçüm gürültüsü (her dakikada az tweet olması), ya da seriyi
\emph{tek bir global trend} (krizden normalleşmeye gidiş) de büyük bir $\Dal$
yaratabilir. Yani büyük $\Dal$ görmek, otomatik olarak ``doğrusal olmayan
dinamik'' demek \emph{değildir}. Bu nedenle iki soruyu ayrı ayrı sormamız ve
\emph{vekil veriler} (surrogate) kullanmamız gerekir:

\begin{itemize}[leftmargin=1.4em]
  \item \textbf{Soru 1:} Büyük $\Dal$ gerçekten \emph{zaman sırasından mı}
        geliyor, yoksa sadece dağılımdan mı? (bkz. \emph{shuffle} testi)
  \item \textbf{Soru 2:} Büyük $\Dal$ sadece \emph{doğrusal korelasyondan} mı
        geliyor, yoksa onun ötesinde bir yapı mı var? (bkz. \emph{IAAFT} testi)
\end{itemize}

Vekil veri, ``gözlenen verinin belirli istatistiksel özelliklerini koruyup
başka bir özelliğini yok eden'' yapay bir seridir. Eğer gözlenen $\Dal$,
vekil verinin üretebileceği $\Dal$ dağılımından anlamlı biçimde büyükse,
``yok edilen'' özellik gerçekten veride vardır. Bu mantık, Bölüm 3.5'te
ayrıntılı kurulacaktır.

% =====================================================================
\section{Yöntem}
% =====================================================================

\subsection{Veri kümesi}
% ---------------------------------------------------------------------

İki bağımsız deprem olayına ait Türkçe tweet kümeleri kullanılmıştır.

\begin{table}[H]
\centering
\caption{Veri kümeleri.}
\label{tab:data}
\begin{tabular}{lrrrr}
\toprule
Olay & Tweet sayısı & Süre & Olay anı (UTC) & Zaman aralığı (UTC)\\
\midrule
Kahramanmaraş (6 Şub 2023) & 942\,180 & 5 gün & 01:17 & 2023-02-06 00:08 -- 02-10 23:59\\
İzmir (30 Eki 2020)        & 971\,295 & 8 gün & 11:51 & 2020-10-29 21:07 -- 11-06 20:59\\
\midrule
\textbf{Toplam} & \textbf{1\,913\,475} & & & \\
\bottomrule
\end{tabular}
\end{table}

Veri kümeleri, Türkçe tweet metni ve zaman damgası içerir (Kaggle kaynaklı,
açık erişimli). İki olayın seçilme gerekçesi: (i) ikisi de çok büyük ve
tetikleyicisi kesin bilinen afetlerdir; (ii) farklı duygusal profil
üretmeleri beklenir (2023: yas/keder; 2020: enkaz ve tsunami kaynaklı korku),
bu da genellenebilirliği güçlendirir.

\subsection{Duygu ölçümü: model seçiminin kritikliği}
% ---------------------------------------------------------------------

\subsubsection{Problem: ``Duygu''yu sayıya çevirmek}
Her tweet'i sayısal duygu değerlerine çevirmek için önceden eğitilmiş
Transformer tabanlı Türkçe duygu/sentiment modelleri kullanılır. Ancak burada
kritik bir sorun vardır: \emph{genel amaçlı bir duygu modeli, afet bağlamında
sistematik olarak yanılabilir.} Örneğin, ``Acil yardım edin!'' ya da
``vefat edenlere rahmet olsun'' cümleleri, felaket bağlamında duygusal
olarak ağır (kaygı, üzüntü) olmasına rağmen, bazı modeller tarafından
sözcük yapısı nedeniyle ``olumlu'' olarak etiketlenebilir.

Bu, saf bir teknik ayrıntı değildir: yanlış bir duygu ölçümü, ardından gelen
\emph{tüm} fraktal analizi geçersiz kılar. Bu nedenle projenin ilk ciddi işi
model doğrulaması olmuştur.

\subsubsection{Üç modelin gerçek veride testi}
Üç aday model, gerçek deprem tweet'leri üzerinde karşılaştırıldı:

\begin{table}[H]
\centering
\caption{Duygu modeli seçimi.}
\label{tab:models}
\begin{tabular}{p{5.2cm}p{6.6cm}p{2.6cm}}
\toprule
Model & Gerçek tweet'lerde gözlenen davranış & Karar\\
\midrule
\texttt{saribasmetehan/...sentiment-analysis} (3 sınıf) &
Felaket bağlamındaki ``yardım/dua/teselli'' dilini yanlışlıkla \emph{pozitif}
okuyor (örn. ``Acil yardım edin!'' $\to +0.98$). &
\textbf{Eleme}\\
\texttt{savasy/bert-base-turkish-sentiment-cased} (2 sınıf) &
Makul: gerçek veride yaklaşık \%70 negatif; yönü doğru. &
\textbf{Robustluk}\\
\texttt{furkankarakuz/turkish-emotion-classifier} (6 duygu) &
Örnek incelemede tutarlı: korku/üzüntü/öfke/sevinç bağlama uygun. &
\textbf{Ana ölçüm}\\
\bottomrule
\end{tabular}
\end{table}

\textbf{Sonuç:} ana duygu ölçümü olarak 6 boyutlu duygu modeli seçildi;
2 sınıflı model ise bağımsız bir \emph{robustluk kontrolü} olarak korundu.
Elenen modelin bulgusu, ``naif model seçimi yanlış bilimsel sonuç üretir''
dersini somutladığı için projenin metodolojik katkılarından biri olarak
raporlanmıştır.

\subsubsection{Kullanılan sinyaller}
6 duygu boyutu (öfke, tiksinme, korku, sevinç, üzüntü, şaşkınlık) doğrudan
kullanılır. Bunlardan iki \emph{bileşik} sinyal türetilir:
\begin{align}
  \text{neg\_emo} &= \text{üzüntü} + \text{korku} + \text{öfke} + \text{tiksinme},\\
  \text{distress} &= \text{üzüntü} + \text{korku},\\
  \text{valence\_emo} &= \text{sevinç} - \text{neg\_emo}.
\end{align}
\textbf{Ana sinyal \texttt{distress}}'tir (üzüntü + korku), çünkü afet
bağlamında en temsilî ve en az gürültülü bileşimdir. Bunların yanında
2 sınıflı modelden \texttt{val\_sav} $= P(\text{poz}) - P(\text{neg})$ ve
salt kontrol amaçlı \texttt{HACİM} $= \log(1+\text{tweet/dk})$ sinyali analiz
edilir.

\subsection{Zaman serisi inşası}
% ---------------------------------------------------------------------

\subsubsection{Dakikalık kutulama}
Her tweet, zaman damgasına göre bir dakikalık ``kutuya'' atanır ve o kutudaki
tüm tweet'lerin duygu puanları \emph{ortalaması} alınır. Böylece dakikada bir
değer içeren düzenli bir seri elde edilir. Kahramanmaraş için $N=7192$,
İzmir için $N=11513$ kutu oluşur. Boş kutular lineer interpolasyonla
doldurulur.

\subsubsection{Neden 1 dakika?}
1 dakika, iki dengeyi gözetir: (i) 5 dakika gibi daha kaba bir çözünürlük kriz
dinamiğinin ince yapısını kaçırır; (ii) 30 saniye gibi daha ince bir çözünürlük
çoğu kutuyu boş bırakır ve ölçüm gürültüsünü artırır. 1 dakika, ``dakikada
onlarca tweet'' yoğunluğuyla ortalama duygu tahmininin güvenilir olduğu en
ince ölçektir.

\subsubsection{Analiz penceresi kuralı (kritik)}
Boş kutuların yoğunlaştığı bölgeler, uzun lineer interpolasyon nedeniyle
\emph{yapay} pürüzsüz parçalar üretir ve bu parçalar büyük ölçeklerde sahte
çoklu fraktallik yaratabilir. İzmir verisinde olay-öncesi gece
(2020-10-29 21:07 $\to$ 10-30 10:00) aşırı seyrektir: 884 dakikanın 586'sı
(yüzde~66) boştur ve bu, ham $\Dal=3.55$ gibi gerçekdışı bir değere yol açar;
seyrek bölge dışlandığında $\Dal \approx 0.67$'ye iner. Bu nedenle her olay
için analiz penceresi
\begin{equation}
  \text{analiz aralığı} = [\,t_{\text{olay}} - 1\text{ saat},\ t_{\text{son}}\,]
\end{equation}
olarak tanımlanmıştır. (Kahramanmaraş'ta gece yeterince yoğun olduğu için bu
kural etkisizdir.)

\subsubsection{Pencereler}
Her olay için dört pencere incelenir:
\begin{itemize}[leftmargin=1.4em]
  \item \textbf{TÜM:} analiz aralığının tamamı;
  \item \textbf{AKUT:} olaydan sonraki ilk 24 saat (kriz);
  \item \textbf{NORM:} serinin son 48 saati (normalleşme);
  \item \textbf{HACİM:} $\log(1+\text{tweet/dk})$ (doğrusal kontrol).
\end{itemize}

\subsection{MF--DFA protokolü ve kalibrasyonu}
% ---------------------------------------------------------------------

\subsubsection{Parametre seçimi}
Kullanılan protokol: $\tau$-derecesi $q \in [-5, 5]$ (21 değer; bakışımlı),
detrend polinom derecesi 2, ölçekler $s_{\min}=32$ ile $s_{\max}=N/4$
arasında logaritmik olarak 30 adet. Bu parametrelerin seçimi rastgele değildir:

\begin{itemize}[leftmargin=1.4em]
  \item \textbf{Ölçek alt sınırı $s_{\min}$:} Gerçek veride $s<24$ dakikada
        güç yasası fitinin doğrusallığı bozuluyor
        ($s_{\min}=16$ iken $R^2=0.63$); $s_{\min}\ge 24$ seçilince
        $R^2\approx 0.94$'e çıkıyor. Bu, küçük ölçekte (kutu-içi örnekleme
        gürültüsünden kaynaklanan) bir ``çapraz-geçiş'' olduğunu gösterir.
        Temiz referans bir çoklu fraktalda ise $R^2=0.99$ ve $\Dal$
        $s_{\min}$'den bağımsızdır. Protokol $s_{\min}=32$ olarak sabitlenmiştir.
  \item \textbf{$q$ aralığı:} $q\in[-5,5]$ ile $q\in[-10,10]$ neredeyse aynı
        $\Dal$'yı verir (doyum gözlenir); bu, ölçümün stabil olduğunu gösterir.
  \item \textbf{Ölçek üst sınırı:} $s_{\max}=N/4$ seçildi (kenar etkilerini
        sınırlamak için).
\end{itemize}

\subsubsection{Bilinen süreçlerle doğrulama}
Motorun doğruluğu üç bağımsız araçla sınandı:
\begin{itemize}[leftmargin=1.4em]
  \item \textbf{Kesirli Gauss gürültüsü (fGn):} Spektral sentezle üretilen ve
        Hurst üssü $H$ bilinen serilerde $h(q{=}2)$ ölçüldü: $H=0.3/0.5/0.7/0.9$
        için sırasıyla $0.317/0.493/0.705/0.894$ bulundu --- teorik değerlerle
        uyumlu (en büyük sapma $0.006$). Aynı serilerin güç spektrumu eğimi,
        ölçülen $-(2H-1)$ ile birebir örtüşür.
  \item \textbf{p-model (binom kaskad):} Teorik $\Dal \approx 1.19$ (p=0.3)
        iken hesaplanan $1.164$; bağıl hata \%2.0. (p=0.2 için teorik $1.996$,
        hesaplanan $1.952$; hata \%2.2.)
  \item \textbf{Multifraktal rastgele yürüyüş (MRW):} Stokastik çoklu fraktal
        referansı olarak üretildi; $\sigma$ arttıkça $\Dal$ monoton artar
        ($0.139/0.547/1.004$, sırasıyla $\sigma=0.3/0.5/0.8$).
\end{itemize}

\subsubsection{Kalibrasyon cetveli}
Aynı protokol, doğası bilinen referans süreçlere uygulanarak bir ``cetvel''
oluşturuldu. Amaç, gerçek verinin $\Dal$'sını tanıdık süreçlerle
karşılaştırmaktır.

\begin{table}[H]
\centering
\caption{Kalibrasyon cetveli ($q\in[-5,5]$, $s_{\min}=32$, $N=8192$). Son
satır gerçek veriden alınan etki büyüklüğüyle karşılaştırma için
eklenmiştir.}
\label{tab:calib}
\begin{tabular}{llccc}
\toprule
Süreç & Karakter & $\Dal$ & $h(q{=}2)$ & IAAFT $z$\\
\midrule
Beyaz gürültü & Tamamen rastgele & 0.076 & 0.515 & $+1.9$\\
Tekdüze dağılım & Tamamen rastgele & 0.016 & 0.491 & ---\\
Rastgele yürüyüş & Belleksiz birikim & 0.095 & 1.439 & ---\\
fGn $H{=}0.3$ & Tekli fraktal & 0.017 & 0.317 & ---\\
fGn $H{=}0.5$ & Tekli fraktal & 0.077 & 0.493 & ---\\
fGn $H{=}0.7$ & Tekli fraktal & 0.025 & 0.705 & ---\\
fGn $H{=}0.9$ & Tekli fraktal & 0.044 & 0.894 & ---\\
AR(1), $\phi{=}0.9$ & Kısa erimli doğrusal hafıza & 0.333 & 0.770 & ---\\
p-model $p{=}0.4$ & Çoklu fraktal (det.) & 0.514 & 0.889 & ---\\
p-model $p{=}0.3$ & Çoklu fraktal (det.) & 1.164 & 0.793 & $+1.0$\\
p-model $p{=}0.2$ & Güçlü çoklu fraktal & 1.952 & 0.656 & ---\\
MRW $\sigma{=}0.8$ & Çoklu fraktal (stok.) & 1.004 & 1.438 & $+3.2$\\
\midrule
\emph{Duygu (distress)} & \emph{gerçek veri} & \emph{0.67--0.78} & --- & \emph{$+10$--$+12$}\\
\bottomrule
\end{tabular}
\end{table}

Bu cetvel üç şeyi birden gösterir. Birincisi, \emph{hiçbir referans sürecin}
ulaşamadığı bir etki büyüklüğü vardır: gerçek verideki $z\approx +10$--$+12$
değerleri, referans süreçlerin en yüksek değeri olan MRW'nin $z=+3.2$'si ile
bile karşılaştırılamaz. İkincisi, $\Dal$ değerleri birbiriyle karıştırılmaz
şekilde ayrışır: rastgele süreçler $\approx 0.1$, tekli fraktallar $\approx 0.1$
altında, gerçek çoklu fraktallar $1$'in üzerindedir; gözlenen $0.67$--$0.78$ değerleri
bu ikisinin arasında, ``orta şiddette çoklu fraktallık'' bölgesindedir. Üçüncüsü
ve en önemlisi, \emph{IAAFT'nin tek başına yeterli bir null olmadığının}
kanıtlanmasıdır (bkz. aşağıdaki kutu).

\begin{quote}
\small
\textbf{Kutu: IAAFT'nin sınırlaması.} p-model kaskadının genlik dağılımı
\emph{ayrıktır} (yalnızca 58 farklı değer) ve IAAFT surrogu bu serinin $\Dal$'sını
neredeyse tamamen taklit eder ($z=+1.0$, anlamsız). Buna karşılık \emph{shuffle}
($z=+4.8$) ve \emph{faz-rastgele} ($z=+20.1$) null'ları çoklu fraktalliği açıkça
yakalar. Bu, bir yöntem hatası değil, IAAFT'in gerçek bir yapısal sınırlamasıdır:
ayrık (kutuplaşmış) genlik dağılımlı serilerde IAAFT null'u gözlemi taklit eder.
Sonuç olarak \textbf{IAAFT tek başına asla yeterli bir null değildir; her zaman
en az bir shuffle veya faz-rastgele null'u ile birlikte raporlanmalıdır.}
Gerçek duygu serilerinin genlik dağılımı düzgün (sürekli) olduğundan, ana
bulgular bu sınırlamadan etkilenmez; nitekim ana bulguda IAAFT $z$'si ile
shuffle $z$'si birbirini doğrulamaktadır.
\end{quote}

\subsection{Surrogate (vekil veri) testleri}
% ---------------------------------------------------------------------

Ortaya çıkan $\Dal$'nın gerçekten önemli olup olmadığını anlamak için bir
``kıyaslama hiyerarşisi'' kuruldu. Her surrogate, orijinal seriden türetilir
ve orijinalin belirli özelliklerini \emph{korur}, belirli özelliklerini
\emph{yok eder}:

\begin{table}[H]
\centering
\caption{Vekil veri (surrogate) hiyerarşisi.}
\label{tab:surr}
\begin{tabular}{lll}
\toprule
Surrogate & Koruduğu & Yok ettiği (test ettiği)\\
\midrule
Shuffle & Genlik dağılımı & Tüm zaman yapısı\\
Faz-rastgele & Güç spektrumu & Faz (doğrusal olmayan yapı)\\
IAAFT & Güç spektrumu + dağılım & Doğrusal olmayan çoklu fraktallik\\
\bottomrule
\end{tabular}
\end{table}

\textbf{Shuffle (karıştırma):} Aynı değerleri rastgele sırayla dizeriz. Böylece
``hangi değerlerin görüldüğü'' korunur, ``hangi sırayla görüldüğü'' yok edilir.
Gözlenen $\Dal$, shuffle null dağılımından anlamlı büyükse, zaman sırası
bilgi taşıyordur.

\textbf{IAAFT} (Iterative Amplitude Adjusted Fourier Transform; Schreiber \&
Schmitz, 1996): Fazları rastgeleleştirir ama hem genlik dağılımını hem de güç
spektrumunu (yani \emph{doğrusal} otokorelasyonu) iteratif olarak korur. Bu
surrogate, ``aynı doğrusal korelasyona sahip, ama doğrusal olmayan yapıdan
yoksun'' bir seriyi temsil eder. Gözlenen $\Dal$, IAAFT null dağılımından
anlamlı büyükse, elimizdeki çoklu fraktallik \emph{doğrusal korelasyonla
açıklanamaz}; yani doğrusal olmayan dinamiklerin izidir.

\subsection{İstatistiksel çıkarım}
% ---------------------------------------------------------------------

Her gerçek pencere için, ilgili surrogate'ten $n=100$ adet üretilir ve her
birinin $\Dal$'sı hesaplanır. Bu, bir \emph{null dağılım} verir. Gözlemi
karşılaştırmak için:
\begin{equation}
  z = \frac{\Dal_{\text{göz}} - \overline{\Dal}_{\text{null}}}
           {\sigma_{\text{null}}}\;,
\end{equation}
ve kritik değerden büyüklerin oranı olarak $p$-değeri. Yorum: büyük $z$
(gözlenen $\Dal$ null'dan çok daha büyük) $\Rightarrow$ gerçek yapı var.
$\Dal$ pozitif tanımlı olduğundan, ``IAAFT ne kadar zayıfsa o kadar iyi'';
$z \approx 0$ ise çoklu fraktallik tamamen doğrusal süreçlerle açıklanabilir
demektir.

\subsection{Kriz ve rutin fazların ayrıştırılması}
% ---------------------------------------------------------------------

Araştırma sorusu 3'ü test etmek için her olay iki faza bölünür:
\begin{itemize}[leftmargin=1.4em]
  \item \textbf{AKUT:} olay anından sonraki ilk 24 saat (kolektif şok).
  \item \textbf{NORM:} serinin son 48 saati (toplumsal normalleşme).
\end{itemize}
Her iki fazda da $\Dal$ ve surrogate testleri ayrı ayrı hesaplanır; aradaki
fark, çoklu fraktalliğin krizle ilişkisini gösterir.

\subsection{Robustluk analizleri}
% ---------------------------------------------------------------------

Büyük $\Dal$'nın dört olası yapay kaynağı sistematik olarak elenmiştir:
\begin{enumerate}[leftmargin=1.6em]
  \item \textbf{Global trend:} Seriden lineer/kuadratik/kübik trend çıkarılır;
        $\Dal$ korunuyorsa gerçek, çöküyorsa trende bağlıdır.
  \item \textbf{Sirkadiyen (24 saatlik) döngü:} İnsan aktivitesi gün içinde
        dalgalandığı için, saat başına ortalaması çıkarılarak test edilir.
  \item \textbf{Ölçek/pencere hassasiyeti:} $q$ aralığı, $s_{\min}$, pencere
        uzunluğu değiştirilerek sonuçların kararlılığı sınanır.
  \item \textbf{Model bağımlılığı:} Farklı duygu modelinden türetilen
        \texttt{val\_sav} sinyaliyle sonuçlar çapraz-doğrulanır.
\end{enumerate}

\subsection{Yazılım, donanım ve tekrarlanabilirlik}
% ---------------------------------------------------------------------

Tüm analiz, açık kaynak Python ile yazılmıştır. MF--DFA motoru, surrogate
üreticileri (IAAFT, shuffle, faz) ve kalibrasyon araçları sıfırdan
kodlanmıştır. Duygu puanlaması PyTorch + Transformers ile GPU üzerinde
yapılmıştır. Sonuçlar, her adımı yeniden üretilebilir kılacak şekilde
kaydedilmiştir (bkz. Ek~\ref{app:yazilim}).

% =====================================================================
\section{Bulgular}
% =====================================================================

\subsection{Çekirdek doğrulama}
% ---------------------------------------------------------------------

Analize geçmeden önce motorun güvenilirliği kanıtlanmıştır. Bilinen süreçlerde
$\Dal$ beklendiği gibi çıkmıştır: beyaz gürültüde neredeyse sıfır, tekli
fraktallarda düşük, p-model'de yüksek (Tablo~\ref{tab:calib}). Bu, ``eğer
gerçek veride büyük $\Dal$ bulursak, bu motorun ürettiği bir artefakt
değildir'' güvencesini verir.

\subsection{Ölçek aralığı ve protokol sağlamlığı}
\label{sec:olcek}
% ---------------------------------------------------------------------

Gerçek duygu serisinde küçük ölçeklerde bir çapraz-geçiş tespit edilmiş ve
analiz penceresi $s_{\min}=32$ ile sınırlandırılmıştır. Bu düzeltme önemlidir:
düzeltilmeden önce $\Dal$ yapay olarak $1.8$ civarında çıkarken, geçerli
ölçek aralığında gerçek değer $\approx 0.7$--$0.9$ mertebesindedir. Global
trend çıkarıldığında sonuç korunur; sirkadiyen çıkarmada da değişmez. Yani
gözlenen yapı ne trend, ne günlük döngü, ne de ölçüm gürültüsüdür.

\subsection{Olay 1: Kahramanmaraş 2023}
% ---------------------------------------------------------------------

Tam seri ve fazlar için $\Dal$ ve surrogate test sonuçları
Tablo~\ref{tab:km} ve Tablo~\ref{tab:kmfaz}'de verilmiştir.

\begin{table}[H]
\centering
\small
\caption{Kahramanmaraş 2023, TÜM penceresi ($N=7192$). $z$ değerleri
surrogate null dağılımına göre standartlaştırılmış sapmadır.}
\label{tab:km}
\begin{tabular}{lccccccc}
\toprule
Sinyal & $\Dal$ & $\Delta h$ & IAAFT $z$ & IAAFT $p$ & Shuffle $z$ & Shuffle $p$ & $N$\\
\midrule
\textbf{distress} (ANA) & 0.779 & 0.502 & $+11.6$ & 0.0099 & $+20.0$ & 0.0099 & 7192\\
üzüntü (sadness)      & 0.414 & 0.195 & $+7.1$  & 0.0099 & $+10.8$ & 0.0099 & 7192\\
korku (fear)           & 0.572 & 0.380 & $+4.5$  & 0.0099 & $+7.1$  & 0.0099 & 7192\\
öfke (anger)           & 1.263 & 0.816 & $+22.6$ & 0.0099 & $+22.3$ & 0.0099 & 7192\\
sevinç (joy)           & 0.873 & 0.543 & $+20.7$ & 0.0099 & $+19.6$ & 0.0099 & 7192\\
valans (val\_sav)      & 0.516 & 0.298 & $+10.4$ & 0.0099 & $+10.1$ & 0.0099 & 7192\\
valans\_emo            & 0.861 & 0.546 & $+18.8$ & 0.0099 & $+20.6$ & 0.0099 & 7192\\
\midrule
HACİM (kontrol)        & 0.997 & 0.700 & $+0.3$  & 0.0693 & $+33.0$ & 0.0099 & 7192\\
\bottomrule
\end{tabular}
\end{table}

\begin{table}[H]
\centering
\small
\caption{Kahramanmaraş 2023, kriz ve normalleşme pencereleri. AKUT: olay
sonrası ilk 24 saat ($N=1432$); NORM: son 48 saat ($N=2881$).}
\label{tab:kmfaz}
\begin{tabular}{lcccc}
\toprule
& \multicolumn{2}{c}{AKUT (24s)} & \multicolumn{2}{c}{NORM (48s)}\\
\cmidrule(lr){2-3}\cmidrule(lr){4-5}
Sinyal & $\Dal$ & IAAFT $z$ & $\Dal$ & IAAFT $z$\\
\midrule
\textbf{distress} & 0.860 & $-0.8$ & 0.378 & $+7.2$\\
üzüntü           & 0.811 & $+0.2$ & 0.422 & $+5.8$\\
korku            & 0.676 & $+1.2$ & 0.289 & $+3.1$\\
öfke             & 1.990 & $+4.3$ & 0.226 & $+3.7$\\
sevinç           & 1.177 & $+4.3$ & 0.114 & $+1.4$\\
val\_sav         & 0.590 & $+0.3$ & 0.266 & $+4.0$\\
valans\_emo      & 0.798 & $+0.5$ & 0.171 & $+2.7$\\
\bottomrule
\end{tabular}
\end{table}

Gözlemler:
\begin{itemize}[leftmargin=1.4em]
  \item Ana sinyal \texttt{distress}, hem \emph{shuffle} hem \emph{IAAFT}
        null'larından anlamlı büyük $\Dal$ verir ($z>10$, $p<0.01$): yani
        çoklu fraktal yapı hem zaman sırasına hem de doğrusal olmayan
        mekanizmalara bağlıdır.
  \item \textbf{Öfke} en yüksek $\Dal$'yı gösterir ($1.26$; akutta $1.99$).
        Bu, ``öfkenin patlamalı/intermittent'' doğasıyla tutarlıdır.
  \item \textbf{Normalleşme fazında} çoklu fraktallik belirgin biçimde
        zayıflar (distress $0.78 \to 0.38$; sevinç $0.87 \to 0.11$). Yani
        kolektif duygu, kriz geçtikçe daha ``rastgele'' bir hâle yaklaşır.
  \item \textbf{AKUT penceresi dikkat gerektirir:} $N=1432$ ile ölçek aralığı
        yalnızca yaklaşık bir dekad olduğundan $\Delta\alpha$ ve $z$ gürültülüdür;
        burada IAAFT $z$'si $-0.8$ çıkmıştır. Bu pencere, bu nedenle
        \emph{betimsel} olarak sunulur; asıl istatistiksel dayanak TÜM ve NORM
        pencereleridir.
  \item \textbf{Kontrol (HACİM):} Tweet hacminin $\Dal$'sı büyük görünse de
        ($\approx 1.0$) IAAFT null'undan anlamlı sapmaz ($z=+0.3$, $p=0.069$).
        Bu, hacmin karmaşıklığının \emph{doğrusal} süreçlerle açıklanabildiğini,
        oysa duygunun karmaşıklığının açıklanamadığını gösterir: duygu,
        yalnızca ``aktivite'' değildir. (Bu kontrolün İzmir olayında daha
        zayıf bir ayrım verdiği ve dolayısıyla \emph{betimsel} bir kontrol olarak
        yorumlandığı bkz. Bölüm~\ref{sec:sinirlilik}.)
\end{itemize}

\subsection{İzmir 2020: İkinci bağımsız olay}
% ---------------------------------------------------------------------

İzmir verisi, aynı protokolle analiz edilmiştir. Duygu profili
Kahramanmaraş'tan \emph{farklıdır}: İzmir'de baskın duygu \textbf{korkudur}
(\%40), üzüntü \%28'dir. Buna karşılık Kahramanmaraş'ta baskın duygu
\textbf{üzüntüdür} (\%37). Bu, iki olayın toplumsal psikolojisinin farklı
olduğunu gösterir (2020: enkaz altında kurtarma ve tsunami korkusu; 2023:
çok geniş çaplı yas).

İzmir verisinde olay-öncesi seyrek gece dışlandığında (bkz.
Bölüm~\ref{sec:olcek}) tüm seri için $\Dal \approx 0.67$ bulunmuştur; bu,
Kahramanmaraş ile \emph{tutarlı} bir mertebedir ve sonuçların tek bir olaya
özgü olmadığını destekler. Tam sonuçlar Tablo~\ref{tab:iz} (TÜM penceresi) ve
Tablo~\ref{tab:izfaz} (kriz ve normalleşme pencereleri) verilmiştir.

\begin{table}[H]
\centering
\small
\caption{İzmir 2020, TÜM penceresi ($N=10689$).}
\label{tab:iz}
\begin{tabular}{lccccccc}
\toprule
Sinyal & $\Dal$ & $\Delta h$ & IAAFT $z$ & IAAFT $p$ & Shuffle $z$ & Shuffle $p$ & $N$\\
\midrule
\textbf{distress} (ANA) & 0.669 & 0.388 & $+10.2$ & 0.0099 & $+21.6$ & 0.0099 & 10689\\
üzüntü (sadness)      & 0.289 & 0.143 & $+4.9$  & 0.0099 & $+7.0$  & 0.0099 & 10689\\
korku (fear)           & 0.515 & 0.289 & $+10.5$ & 0.0099 & $+16.2$ & 0.0099 & 10689\\
öfke (anger)           & 0.666 & 0.391 & $+19.6$ & 0.0099 & $+13.2$ & 0.0099 & 10689\\
sevinç (joy)           & 0.841 & 0.514 & $+10.4$ & 0.0099 & $+17.5$ & 0.0099 & 10689\\
valans (val\_sav)      & 0.428 & 0.200 & $+12.6$ & 0.0099 & $+9.8$  & 0.0099 & 10689\\
valans\_emo            & 0.797 & 0.469 & $+20.5$ & 0.0099 & $+20.7$ & 0.0099 & 10689\\
\midrule
HACİM (kontrol)        & 1.019 & 0.587 & $+9.7$  & 0.0099 & $+40.4$ & 0.0099 & 10689\\
\bottomrule
\end{tabular}
\end{table}

\begin{table}[H]
\centering
\small
\caption{İzmir 2020, kriz ve normalleşme pencereleri. AKUT: olay sonrası ilk
24 saat ($N=1440$); NORM: son 48 saat ($N=2881$).}
\label{tab:izfaz}
\begin{tabular}{lcccc}
\toprule
& \multicolumn{2}{c}{AKUT (24s)} & \multicolumn{2}{c}{NORM (48s)}\\
\cmidrule(lr){2-3}\cmidrule(lr){4-5}
Sinyal & $\Dal$ & IAAFT $z$ & $\Dal$ & IAAFT $z$\\
\midrule
\textbf{distress} & 0.745 & $+6.7$ & 0.246 & $+2.7$\\
üzüntü           & 0.284 & $+1.7$ & 0.282 & $+3.7$\\
korku            & 0.306 & $+0.4$ & 0.130 & $+1.2$\\
öfke             & 0.642 & $+6.5$ & 0.324 & $+3.3$\\
sevinç           & 0.919 & $+8.1$ & 0.159 & $+0.7$\\
val\_sav         & 0.600 & $+8.4$ & 0.177 & $+1.8$\\
valans\_emo      & 1.010 & $+10.7$ & 0.168 & $+0.5$\\
\bottomrule
\end{tabular}
\end{table}

İzmir'de de aynı nitel desen görülür: \emph{tüm} sinyaller TÜM penceresinde
IAAFT ve shuffle null'larını anlamlı biçimde aşar; en güçlü sinyal yine
öfke/valans\_emo'dur ($z\approx +20$); normalleşme fazında $\Dal$ düşer
(distress $0.67 \to 0.25$) ve $z$'ler büyük ölçüde anlamsızlaşır. Kriz
penceresinde ise İzmir'de $z$'ler Kahramanmaraş'a göre \emph{daha yüksek}
($+6.7$ vs. $-0.8$); bunun nedeni iki olayın farklı duygusal profilleri
(İzmir'de korku ve sevinç baskın, Kahramanmaraş'ta üzüntü ve öfke) ve
dolayısıyla farklı interim (aralıksız) yapılarıdır.

Yalnızca bir fark önemlidir: İzmir'de \textbf{hacim kontrolü de} IAAFT null'unu
aşmaktadır ($z=+9.7$). Bu, ``hacim doğrusaldır'' argümanının İzmir'de
geçerli olmadığını gösterir; bu yüzden hacim kontrolü raporun tamamında
\emph{betimsel} bir karşılaştırma olarak sunulmuş, ana argüman ise doğrudan
IAAFT ve shuffle null'larının \emph{çift} karşılaştırmasına dayandırılmıştır
(bkz. Bölüm~\ref{sec:sinirlilik}).

\subsection{Olaylar arası karşılaştırma ve doğrusal/doğrusal-olmayan ayrımı}
% ---------------------------------------------------------------------

Her iki olayda da ortak desen şudur: duygu sinyalleri TÜM penceresinde hem
IAAFT hem shuffle null'larını anlamlı biçimde aşar, ve normalleşme fazında bu
sapma kapanır. Karşılaştırmanın özeti Tablo~\ref{tab:kars}'te verilmiştir.

\begin{table}[H]
\centering
\small
\caption{İki olayın karşılaştırması (ana sinyal \texttt{distress}).}
\label{tab:kars}
\begin{tabular}{lcccc}
\toprule
& \multicolumn{2}{c}{Kahramanmaraş 2023} & \multicolumn{2}{c}{İzmir 2020}\\
\cmidrule(lr){2-3}\cmidrule(lr){4-5}
Pencere & $\Dal$ & IAAFT $z$ & $\Dal$ & IAAFT $z$\\
\midrule
TÜM  & 0.779 & $+11.6$ & 0.669 & $+10.2$\\
AKUT & 0.860 & $-0.8$  & 0.745 & $+6.7$\\
NORM & 0.378 & $+7.2$  & 0.246 & $+2.7$\\
HACİM& 0.997 & $+0.3$  & 1.019 & $+9.7$\\
\bottomrule
\end{tabular}
\end{table}

Tablo~\ref{tab:kars} ve Tablo~\ref{tab:surr} birlikte şu sonuca götürür:
\emph{kolektif duygunun karmaşıklığı, bir doğrusal gürültü sürecine
indirgenemez.} İki olayın $z$'leri birbirine çok yakındır ($+11.6$ ve $+10.2$);
bu, farklı yıl, farklı ölçek, farklı duygusal profil ve farklı tweet sayısına
sahip iki bağımsız olayda aynı nicel yapının ortaya çıktığını gösterir.

Hacim kontrolü hakkında bir ayrıntı gerekir. Kahramanmaraş'ta hacim
$\log(1+\text{tweet/dk})$ serisi IAAFT ile açıklanırken ($z=+0.3$), İzmir'de
aynı kontrol anlamlı biçimde sapar ($z=+9.7$). İncelendiğinde bunun nedeni
Kahramanmaraş serisindeki 60 adet sıfır-dakikanın IAAFT null'unu
kararsızlaştırması olduğu görülmüştür ($z$ aynı veriyle $40$ surrogal
arasında $+0.3$ ile $+1.3$ arasında değişmektedir). Bu nedenle raporun ana
argümanı ``duygu doğrusaldır, hacim değil'' biçiminde kurulmamış; bunun yerine
duygunun kendi null modellerine karşı test edilmesine dayandırılmıştır.

\subsection{Robustluk bulguları}
% ---------------------------------------------------------------------

\begin{itemize}[leftmargin=1.4em]
  \item Trend (lineer/kuadratik/kübik) çıkarıldığında $\Dal$ korunmuştur:
        sonuç global trendin ürünü değildir.
  \item Sirkadiyen döngü çıkarıldığında $\Dal$ değişmemiştir: 24 saatlik
        aktivite ritmi bir artefakt kaynağı değildir.
  \item $q$ aralığı genişletildiğinde $\Dal$ doyuma ulaşır: ölçüm stabildir.
  \item İkinci duygu modelinden (\texttt{val\_sav}) elde edilen seri,
        benzer yönde sonuç verir: ölçüm modelden bağımsızdır.
\end{itemize}

% =====================================================================
\section{Tartışma}
% =====================================================================

\subsection{Ana bulgu: Duygu serisi çoklu fraktaldır}
% ---------------------------------------------------------------------

Elde edilen sonuçlar, araştırma sorularının çoğunu olumlu yanıtlar. Kolektif
duygu zaman serisi (i) saf rastgele \emph{değildir}, (ii) sadece doğrusal
korelasyonla da açıklanamaz, (iii) kriz ve normalleşme fazlarında farklı
davranır, (iv) bu nitel desen iki bağımsız olayda ve iki farklı duygu
modeliyle tekrarlanır. Bu dört madde birlikte, H1--H4 hipotezlerini
destekler.

Bunun düşünsel anlamı şudur: sosyal medya duygusu, ``milyonlarca bağımsız
kişinin bağımsız kararları'' olarak indirgenebilecek bir gürültü değildir;
içinde \emph{doğrusal olmayan, ölçekten bağımsız, uzun erimli bir düzen}
barındırır. Bu, kalabalık insan davranışlarının bazı yönleriyle türbülans ve
finansal piyasa gibi ``kendine örgütlenen'' sistemlere benzediğini gösterir.

\subsection{Kriz ve normalleşme}
% ---------------------------------------------------------------------

Çoklu fraktalliğin kriz anında güçlenip normalleşme sürecinde zayıflaması,
\emph{kolektif duygunun bir ``karmaşıklık imzası'' olduğunu} ve bu imzanın
toplumsal iyileşme süreciyle birlikte çözüldüğünü göstermektedir. Bu,
gelecekte bir ``erken uyarı göstergesi'' (toplumsal gerginlik/panik
yoğunlaşması) geliştirmeye kapı aralar: $\Dal$'daki ani yükselişler, olağan
dışı bir toplumsal hareketlenmenin işareti olabilir.

\subsection{Duygu, aktivite yoğunluğundan farklıdır}
% ---------------------------------------------------------------------

Belki de en \emph{pratik} bulgu budur: İçerik (duygu) ile miktar (tweet
sayısı) ölçeklenebilir birbirinden farklı yapılar taşımaktadır. Tweet hacmi
serisi de çoklu fraktaldir ($\Delta\alpha\approx 1.0$, yani duygu
serilerininkinden \emph{daha} geniş bir spektrum), fakat bu karmaşıklığın
doğrusal (otokorelasyonlu) bir gürültü sürecine indirgenip indirgenemediği
olaylar arasında değişmektedir. Dolayısıyla ``hacim de çoklu fraktaldır''
gözlemi tek başına ``duygu da öyle'' gözlemini doğrulamaz.

Bununla birlikte, duygu serilerinin iki bağımsız olayda da hem IAAFT hem de
shuffle null'larını $z\approx +10$--$+22$ ile aşması, yani \emph{kendi}
null modelleri karşısında açık biçimde sapması, asıl istatistiksel ayrımı
verir. Bu ayrım, ``sadece tweet sayısına bakarak duygu analizi yapmak''
yaklaşımının yetersiz olacağını gösteresi bakımından değerlidir. Yine de
bu tür bir karşılaştırmanın, hacim kontrolünün iki olayda tutarsız davrandığı
göz önüne alınarak \emph{eşit ölçekte bir kanıt} olarak değil, \emph{yardımcı
bir gözlem} olarak sunulması gerektiğini vurgulamak isteriz.

\subsection{Metodolojik dersler}
% ---------------------------------------------------------------------

\begin{enumerate}[leftmargin=1.6em]
  \item \textbf{Model seçimi kritiktir.} Felaket bağlamında sosyal medya
        duygu ölçümünde, ``yardım/dua/teselli'' dilini yanlış okuyan naif bir
        model, tüm fraktal analizi geçersiz kılabilir. Bu çalışma, üç modeli
        gerçek veride karşılaştırarak bu tuzağı somutlamıştır.
  \item \textbf{Ölçek aralığı $R^2$ ile sabitlenmelidir.} Küçük ölçekteki
        çapraz-geçişi görmezden gelmek, $\Dal$'yı olduğundan büyük gösterir.
  \item \textbf{Seyrek veri bölgeleri dışlanmalıdır.} Uzun interpolasyonlar
        yapay çoklu fraktallik üretir; bu, $\Dal$'yı $5$ katına kadar
        şişirebilir.
  \item \textbf{Null modelleri çoğaltmak gerekir.} IAAFT, ayrık genlik
        dağılımlı serilerde çoklu fraktalliği yok edemeyebilir (bu çalışmada
        p-model üzerinde belgelendi: IAAFT $z=+1.0$, shuffle $z=+4.8$,
        faz-rastgele $z=+20.1$). Bu nedenle tek bir null modeline güvenmek
        yerine en az iki bağımsız null (burada IAAFT ve shuffle) birlikte
        raporlanmalıdır. Bu ilke, projenin genel olarak kullanılabilecek
        metodolojik katkılarından biridir.
  \item \textbf{Bir kontrolün ``anlamlı sapmaması'' tek başına yeterlidir.}
        Buradaki hacim kontrolü, iki olay arasında farklı davranmış ve
        ``duygu doğrusaldır'' biçimindeki basit bir argüman bu yüzden
        kurulmamıştır. Kontroller, eşit ölçekte bir kanıt olarak değil,
        yardımcı gözlem olarak sunulmalıdır.
\end{enumerate}

\subsection{Sınırlılıklar}
\label{sec:sinirlilik}
% ---------------------------------------------------------------------

\begin{itemize}[leftmargin=1.4em]
  \item \textbf{Duygu ölçümü tam değildir:} Modeller bağlamı sınırlı anlar.
        Bunu azaltmak için çapraz-model doğrulaması yapılmıştır.
  \item \textbf{Temsilî olmayan örneklem:} Sosyal medya kullanıcıları,
        nüfusun tamamını değil, kendini ifade eden bir kesimini temsil eder.
  \item \textbf{İki olayla sınırlı genellenebilirlik:} Sonuçlar iki olayla
        tutarlıdır, ancak daha geniş bir olay setiyle pekiştirilmelidir.
  \item \textbf{Nonstationarity:} Kriz$\to$normalleşme geçişi büyük ölçekleri
        etkiler; bu, detrend analizleriyle kontrol edilmiştir.
  \item \textbf{IAAFT'nin ayrık genlik dağılımındaki sınırlaması:} IAAFT
        null'u, kutuplaşmış (ayrık) genlik dağılımlı serilerde gözlenen
        çoklu fraktalliği taklit edebilmektedir (p-model üzerinde belgelendi).
        Bu çalışmada tüm sonuçlar bu nedenle \emph{çift} null testiyle
        (IAAFT + shuffle) verilmiştir; tek başına IAAFT'a dayanılmamıştır.
  \item \textbf{Hacim kontrolünün tutarsızlığı:} HACİM serisi iki olayda
        farklı davranmış (Kahramanmaraş $z=+0.3$, İzmir $z=+9.7$); bu
        kontrol yalnızca betimsel olarak yorumlanabilmiştir.
  \item \textbf{AKUT penceresinin kısalığı:} $N\approx 1440$ ile ölçek
        aralığı yaklaşık bir dekadla sınırlı olduğundan, bu penceredeki
        $\Delta\alpha$ ve $z$ değerleri gürültülüdür; AKUT penceresi
        betimsel bulgu olarak sunulmuştur.
  \item \textbf{Duygu ölçümünün hacimle ortak değişkenliği:} Aynı tweet'lerden
        türetildiği için duygu serisinin dakika-içi varyansı, tweet yoğunluğu
        ile ortak değişkenliğe sahiptir. IAAFT ve shuffle null'ları bu ortak
        değişkenliğin büyük kısmını elese de, tam olarak ortadan kaldırdığı
        söylenemez. Daha kesin bir kontrol, hacim-konumlandırılmış bir null
        modeldir (gelecek çalışma).
\end{itemize}

\subsection{Genellenebilirlik}
% ---------------------------------------------------------------------

Farklı yıl, büyüklük ve duygusal profile sahip iki bağımsız olayda aynı
nitel desenin (anlamlı $\Dal$, krizde güçlenme) görülmesi, bulgunun bir
olaya özgü tesadüf olmadığını, kolektif duygu dinamiğinin genel bir özelliği
olabileceğini işaret eder.

% =====================================================================
\section{Sonuç ve Gelecek Çalışmalar}
% =====================================================================

Bu proje, basit ama derin bir soruyla başladı: ``Bir deprem sırasında
milyonlarca insanın birlikte ürettiği duygu, rastgele midir?'' Yanıt, hayır
çıktı. Kolektif duygu zaman serisi, hem zaman sırasına hem de doğrusal
olmayan mekanizmalara bağlı, ölçekten bağımsız, çoklu fraktal bir yapı taşır;
bu yapı kriz anında güçlenir, normalleşmeyle zayıflar ve iki bağımsız olayda
tekrarlanır.

Bu çalışmanın yöntem açısından da bir katkısı vardır: (i) felaket bağlamında
duygu ölçümünde naif model seçiminin yarattığı ölçüm hatasının somut olarak
gösterilmesi, (ii) ayrık genlik dağılımlı serilerde IAAFT null'unun yetersiz
kaldığının deneysel olarak belgelenmesi ve buna karşılık çift null
(IAAFT + shuffle) kullanımının önerilmesi, (iii) seyrek veri bölgelerinin
yapay çoklu fraktallik üretmesi nedeniyle analiz penceresi kuralının
sabitlenmesi.

\textbf{Gelecek çalışmalar:}
\begin{enumerate}[leftmargin=1.6em]
  \item \textbf{Daha çok olay:} Farklı türde krizler (yangın, sel, salgın)
        eklenerek genellenebilirlik genişletilebilir.
  \item \textbf{Erken uyarı göstergesi:} $\Dal$'daki ani yükseliş, gerçek
        zamanlı toplumsal gerginlik uyarısı olarak test edilebilir.
  \item \textbf{İleri yöntemler:} A-MFDFA (asalama etkilerini ayırma),
        çok değişkenli MF--DFA, ağ analiziyle birlikte kullanım.
  \item \textbf{Coğrafi/yapısal ayrıştırma:} Kurtarma, kayıp, haber, destek
        gibi alt-temalara göre seriler bölünebilir.
\end{enumerate}

% =====================================================================
\section*{Kaynaklar}
\addcontentsline{toc}{section}{Kaynaklar}
% =====================================================================
\begin{thebibliography}{99}
\bibitem{kantelhardt2002} J. W. Kantelhardt, S. A. Zschiegner, E. Koscielny-Bunde,
A. Havlin, S. Bunde, H. E. Stanley, ``Multifractal detrended fluctuation
analysis of nonstationary time series,'' \emph{Physica A} \textbf{316}, 87--114, 2002.
\bibitem{peng1994} C.-K. Peng, S. V. Buldyrev, S. Havlin, M. Simons,
H. E. Stanley, A. L. Goldberger, ``Mosaic organization of DNA nucleotides,''
\emph{Physical Review E} \textbf{49}, 1685, 1994.
\bibitem{peng1995} C.-K. Peng et al., ``Long-range correlations in nucleotide
sequences'' ve ``Quantification of scaling exponents ... heartbeat,'' 1995.
\bibitem{schreiber1996} T. Schreiber, A. Schmitz, ``Improved surrogate data for
nonlinearity tests,'' \emph{Physical Review Letters} \textbf{77}, 635, 1996.
\bibitem{mandelbrot1982} B. B. Mandelbrot, \emph{The Fractal Geometry of Nature},
W. H. Freeman, 1982.
\bibitem{meneveau1987} C. Meneveau, K. R. Sreenivasan, ``Simple multifractal
cascade model for the dissipation rate,'' \emph{Physical Review A} \textbf{36}, 1987.
\bibitem{bacry2001} E. Bacry, J. Delour, J.-F. Muzy, ``Multifractal random
walk,'' \emph{Physical Review E} \textbf{64}, 026103, 2001.
\bibitem{calvet2004} L. Calvet, A. Fisher, ``How to forecast long-run
volatility,'' \emph{Journal of Financial Econometrics}, 2004.
\bibitem{ihlen2012} E. A. F. Ihlen, ``Introduction to multifractal detrended
fluctuation analysis in Matlab,'' \emph{Frontiers in Physiology} \textbf{3}, 141, 2012.
\bibitem{stanley1999} H. E. Stanley et al., ``Statistical physics and
physiology: from DNA to heartbeats,'' 1999.
\bibitem{kantelhardt2002b} D. Rybski, S. V. Buldyrev, S. Havlin, F. Liljeros,
H. A. Makse, ``Clustering of human interaction'' (sosyal sistem örneği).
\bibitem{kramer2014} A. D. I. Kramer, J. E. Guillory, J. T. Hancock,
``Experimental evidence of massive-scale emotional contagion through social
networks,'' \emph{PNAS} \textbf{111}, 2014.
\bibitem{golder2011} S. A. Golder, M. W. Macy, ``Diurnal and seasonal mood
vary with work, sleep, and daylength across diverse cultures,''
\emph{Science} \textbf{333}, 2011.
\end{thebibliography}

% =====================================================================
\appendix
\section{Formüllerin ayrıntılı türetimi}
\label{app:formul}
% =====================================================================

\textbf{(a) $q=0$ dalının türetimi.} \eqref{eq:Fq} tanımında $q\to 0$
limitini alalım. $x^{q} = e^{q\ln x} \approx 1 + q\ln x$ olduğundan,
\begin{equation}
  \frac{1}{q}\ln\!\left[\frac{1}{N_s}\sum_\nu (F^2)^{q/2}\right]
  \xrightarrow{q\to 0} \frac{1}{2N_s}\sum_\nu \ln F^2(\nu,s),
\end{equation}
bu da $\log$-ortalamanın üstelidir. Bu yüzden $q=0$ dalı
$F_0(s)=\exp\!\bigl[\tfrac{1}{2N_s}\sum \ln F^2\bigr]$ olarak yazılır.
(Bu projede bu dal bağımsız doğrulamayla düzeltilmiştir; önceki hatalı sürüm
$h(0)$'ı iki katı hesaplıyordu.)

\vspace{0.5em}
\textbf{p-model genişliği.} Binom kaskadında her adımda kütle $p$ ve $1-p$
oranında ikiye bölünür. Ölçek üssü
$\tau(q) = -\ln\!\bigl(p^q + (1-p)^q\bigr)/\ln 2$ olur; Legendre dönüşümüyle
$\alpha_{\min} = -\log_2 p$ ve $\alpha_{\max} = -\log_2(1-p)$ elde edilir.
Böylece
\begin{equation}
  \Dal = \alpha_{\max} - \alpha_{\min} = \log_2\frac{1-p}{p}.
\end{equation}
$p=0.3$ için $\Dal \approx 1.19$; bu, hesaplanan $1.164$ ile uyuşur
(bağıl hata \%2.0).

\section{Yazılım mimarisi}
\label{app:yazilim}
\begin{itemize}[leftmargin=1.4em]
  \item \texttt{run\_all.py} --- \emph{tek giriş noktası}: veri denetimi,
        skorlama, seri üretimi, MF--DFA analizi, kalibrasyon, karşılaştırma ve
        testleri sırayla çalıştıran uçtan uca boru hattı.
  \item \texttt{mfdfa\_core.py} --- MF--DFA motoru, surrogate üreticileri
        (IAAFT, shuffle, faz-rastgele) ve referans süreç üreticileri
        (fGn, AR(1), beyaz gürültü, MRW, p-model).
  \item \texttt{calibration.py}, \texttt{verify\_fix.py} --- doğrulama ve
        kalibrasyon cetveli.
  \item \texttt{score2.py} --- duygu puanlaması (iki olay, iki model).
  \item \texttt{build\_series2.py} --- dakikalık zaman serisi üretimi.
  \item \texttt{analyze2.py}, \texttt{run\_analysis.py} --- MF--DFA + surrogate
        analiz sürücüsü.
  \item \texttt{compare\_events2.py} --- olaylar arası karşılaştırma.
  \item \texttt{tests/test\_mfdfa.py} --- 28 adet birim testi: Hurst kurtarımı,
        p-model teorik karşılaştırması, $q=0$ dalı, surrogate testleri,
        IAAFT sınırlamasının belgelenmesi ve kayıtlı sonuçların bütünlüğü.
  \item \texttt{requirements.txt} --- sabitlenmiş bağımlılık sürümleri.
  \item Robustluk: \texttt{robust\_check.py}, \texttt{scale\_test.py},
        \texttt{r2\_check.py}, \texttt{diurnal\_check.py}, \texttt{qsensitivity.py}.
\end{itemize}

\end{document}
